\documentclass[12pt,letterpaper]{article}
\usepackage[utf8]{inputenc}
\usepackage[spanish]{babel}
\usepackage{graphicx}
\usepackage{listings}
\usepackage{hyperref}
\usepackage{amsmath}
\usepackage{geometry}
\usepackage{booktabs}
\usepackage{xcolor}

% Configuración de márgenes
\geometry{left=2.5cm,right=2.5cm,top=3cm,bottom=3cm}

% Configuración de colores para código (listings)
\definecolor{codegreen}{rgb}{0,0.6,0}
\definecolor{codegray}{rgb}{0.5,0.5,0.5}
\definecolor{codepurple}{rgb}{0.58,0,0.82}
\definecolor{backcolour}{rgb}{0.95,0.95,0.92}

\lstdefinestyle{pythonstyle}{
    backgroundcolor=\color{backcolour},   
    commentstyle=\color{codegreen},
    keywordstyle=\color{magenta},
    numberstyle=\tiny\color{codegray},
    stringstyle=\color{codepurple},
    basicstyle=\ttfamily\footnotesize,
    breakatwhitespace=false,         
    breaklines=true,                 
    captionpos=b,                    
    keepspaces=true,                 
    numbers=left,                    
    numbersep=5pt,                  
    showspaces=false,                
    showstringspaces=false,
    showtabs=false,                  
    tabsize=2
}
\lstset{style=pythonstyle}

\title{
    \textbf{Informe Técnico: Análisis y Diseño de Resiliencia}\\
    \large Práctica de Tolerancia a Fallas en Microservicios
}
\author{
    \textbf{Integrante 1:} Nombre Apellido \\
    \textbf{Integrante 2:} Nombre Apellido \\
    \vspace{0.5cm}\\
    Curso de Sistemas Distribuidos \\
    Universidad/Institución
}
\date{\today}

\begin{document}

\maketitle
\thispagestyle{empty}
\newpage

\tableofcontents
\newpage

\section{Introducción}
Este informe técnico aborda el análisis teórico y diseño de soluciones a nivel de producción para los dos escenarios de fallo que no fueron implementados en el código práctico de la arquitectura del clúster. 

Los fallos analizados bajo fundamentos de sistemas distribuidos y resiliencia en este documento son:
\begin{itemize}
    \item \textbf{Fallo 4: Base de Datos Intermitente (Conectividad)}
    \item \textbf{Fallo 6: Condición de Carrera (Consistencia)}
\end{itemize}

---

\section{Análisis de Fallo 4: Base de Datos Intermitente (Conectividad)}

\subsection{Explicación Teórica y Fundamentos Distribuidos}
% Pauta de redacción:
% Explique el fallo utilizando conceptos teóricos de conectividad y redes.
% Detalle cómo las intermitencias físicas de red (pérdida de paquetes, latencia transitoria, fallos de handshake) 
% o la saturación del pool de conexiones del socket de base de datos impiden la persistencia.
% Relacione este comportamiento con los conceptos de fallos transitorios vs permanentes.

Escriba aquí la explicación teórica...

\subsection{Relación con Clases e Investigación MLR}
% Pauta de redacción:
% Relacione las intermitencias con las lecturas académicas del MLR o las diapositivas de clase.
% Explique la relevancia del manejo de fallos transitorios mediante patrones defensivos.

Escriba aquí la relación con la literatura...

\subsection{Solución Propuesta a Nivel de Producción}
% Pauta de redacción:
% Proponga una arquitectura o solución técnica detallada de nivel producción.
% Sugerencias:
% 1. Mecanismo de Reintentos con Backoff Exponencial y Jitter (evitar el efecto "Thundering Herd").
% 2. Implementación del Patrón Transactional Outbox para desacoplar escrituras en base de datos.
% 3. Uso de un Connection Pooler (como PgBouncer para PostgreSQL) para mitigar la sobrecarga de conexiones.
% 4. Idempotencia en API mediante llaves de idempotencia (Idempotency Keys).

Escriba aquí la solución de producción propuesta...

\subsection{Pseudocódigo o Diagrama de Solución}
A continuación se presenta el pseudocódigo que ilustra la implementación de reintentos con backoff exponencial y jitter en la conexión de base de datos:

\begin{lstlisting}[language=Python, caption=Pseudocódigo de reconexión y reintentos con Jitter]
import time
import random

def conectar_con_backoff(max_intentos=5, base_delay=1.0, max_delay=30.0):
    intento = 0
    while intento < max_intentos:
        try:
            # Intentar conexion a PostgreSQL
            db = conectar_postgres()
            return db
        except Exception as e:
            intento += 1
            if intento == max_intentos:
                raise e
            # Calcular delay exponencial con jitter aleatorio
            temp = min(max_delay, base_delay * (2 ** intento))
            sleep_time = temp / 2 + random.uniform(0, temp / 2)
            print(f"Error de conexion. Reintentando en {sleep_time:.2f} segundos...")
            time.sleep(sleep_time)
\end{lstlisting}

% Para incluir un diagrama visual en formato imagen, descomente las siguientes lineas:
% \begin{figure}[h!]
%     \centering
%     \includegraphics[width=0.8\textwidth]{diagramas/solucion_fallo4.png}
%     \caption{Flujograma del patrón Transactional Outbox para resiliencia de datos.}
%     \label{fig:outbox}
% \end{figure}

---

\section{Análisis de Fallo 6: Condición de Carrera (Consistencia)}

\subsection{Explicación Teórica y Fundamentos Distribuidos}
% Pauta de redacción:
% Explique por qué se genera una condición de carrera cuando múltiples peticiones leen y modifican el mismo recurso (asiento).
% Relacione este fallo directamente con el Teorema CAP (Consistencia, Disponibilidad, Tolerancia a Particiones).
% Explique cómo la falta de sincronización distribuye inconsistencias (ej: doble venta de boleto).

Escriba aquí la explicación teórica...

\subsection{Relación con Clases e Investigación MLR}
% Pauta de redacción:
% Relacione la inconsistencia y los modelos de aislamiento de transacciones SQL (Read Committed, Serializable) 
% con lo visto en las clases teóricas de concurrencia distribuida.

Escriba aquí la relación con la literatura...

\subsection{Solución Propuesta a Nivel de Producción}
% Pauta de redacción:
% Describa la solución ideal para producción.
% Sugerencias:
% 1. Bloqueo Pesimista (Pessimistic Locking): SELECT ... FOR UPDATE en la fila del asiento en SQL.
% 2. Patrón de Reserva Temporal (Holds) con expiración (ej: Redis TTL de 10 minutos).
% 3. Bloqueo Optimista (Optimistic Locking) mediante versionamiento de registros.

Escriba aquí la solución de producción propuesta...

\subsection{Pseudocódigo o Diagrama de Solución}
A continuación se ilustra mediante pseudocódigo la solución utilizando bloqueo pesimista a nivel de base de datos para asegurar el procesamiento atómico de la reserva:

\begin{lstlisting}[language=Python, caption=Pseudocódigo de control de concurrencia con Bloqueo Pesimista]
# Flujo atomico de reserva en base de datos PostgreSQL
def procesar_reserva_concurrente(db_connection, asiento_id, usuario_id):
    with db_connection.cursor() as cursor:
        # 1. Bloquear la fila del asiento especifico (FOR UPDATE)
        cursor.execute(
            "SELECT disponible FROM asientos WHERE id = %s FOR UPDATE", 
            (asiento_id,)
        )
        row = cursor.fetchone()
        
        if row and row[0] is True: # Si esta disponible
            # 2. Descontar del inventario y registrar la reserva
            cursor.execute(
                "UPDATE asientos SET disponible = FALSE WHERE id = %s", 
                (asiento_id,)
            )
            cursor.execute(
                "INSERT INTO reservas (id, asiento_id, usuario_id, estado) VALUES (%s, %s, %s, 'CONFIRMADA')",
                (generar_uuid(), asiento_id, usuario_id)
            )
            db_connection.commit()
            return {"status": "Reserva exitosa"}
        else:
            db_connection.rollback()
            return {"status": "Fallo: Asiento ya reservado"}
\end{lstlisting}

% Para incluir un diagrama de secuencia de concurrencia en formato imagen, descomente las siguientes lineas:
% \begin{figure}[h!]
%     \centering
%     \includegraphics[width=0.8\textwidth]{diagramas/solucion_fallo6.png}
%     \caption{Diagrama de secuencia de la reserva atómica con bloqueo pesimista.}
%     \label{fig:concurrencia}
% \end{figure}

---

\section{Conclusiones}
% Pauta de redacción:
% Redacte las conclusiones generales del análisis y su importancia en sistemas de reservas de alta concurrencia.

Escriba aquí las conclusiones...

---

\section{Referencias Bibliográficas}
% Pauta de redacción:
% Cite la literatura académica de su investigación MLR y diapositivas de clase correspondientes.
\begin{thebibliography}{9}
    \bibitem{cap} 
    Fox, A., \& Brewer, E. A. (1999). \textit{Harvest, yield, and scalable tolerant systems}. Proceedings of the Seventh Workshop on Hot Topics in Operating Systems.
    
    \bibitem{mlr} 
    Referencia a tu documento MLR...
    
    \bibitem{clase}
    Diapositivas de clase: Tema de Tolerancia a Fallas y Concurrencia Distribuida.
\end{thebibliography}

\end{document}
